\documentclass[../../main.tex]{subfiles}
\begin{document}
% SOURCE p=248; printed=233; object=ch07-notes-products-references; region=top
\paragraph*{\cref{section:ch07-radon-products}的注记。}
这里给出的 Radon 乘积 Fubini--Tonelli 定理实质上归功于 deLeeuw\cite{folland-bib-031}；另见 Cohn\cite[\S7.6]{folland-bib-027}。Hewitt 与 Ross\cite[\S13]{folland-bib-075}给出了另一版本，其中还包含非 $\sigma$ 有限情形的一些进一步结果。
% SOURCE p=248; printed=233; object=ch07-notes-general-probability-products; region=upper-middle
\cref{theorem:ch07-product-regular-borel}实质上归功于 Nelson\cite{folland-bib-103}。这个结果也有纯测度论版本：若 $\{(X_\alpha,\mathscr M_\alpha,\mu_\alpha)\}_{\alpha\in A}$ 为测度空间族，且每个 $\mu_\alpha(X_\alpha)=1$，则可以在 $(\prod_\alpha X_\alpha,\bigotimes_\alpha\mathscr M_\alpha)$ 上定义乘积测度 $\prod_\alpha\mu_\alpha$。参见 Saeki\cite{folland-bib-127}、Halmos\cite[\S38]{folland-bib-061}或 Hewitt 与 Stromberg\cite[\S22]{folland-bib-076}。
% SUPPLEMENT: id=v14-ch07-notes-general-products-proof; parent=ch07-notes-general-probability-products; kind=omitted-proof+motivation+hypotheses; origin=independent-derivation
\begin{remark}
\textbf{任意概率空间族的乘积构造（重构新增）。}对有限坐标集 $B\subseteq A$，记 $X_B=\prod_{\alpha\in B}X_\alpha$、$\mathscr M_B=\bigotimes_{\alpha\in B}\mathscr M_\alpha$，以 $\mu_B$ 表示有限乘积概率测度。空坐标乘积为单点，质量为 $1$。所有柱集 $\pi_B^{-1}(E)$（$E\in\mathscr M_B$）组成一个代数 $\mathcal C$：补集保留同一坐标，有限并可提升到坐标并集后计算。定义
\begin{align*}
P(\pi_B^{-1}(E))=\mu_B(E).
\end{align*}
若同一柱集有两种表示，把它们提升到两坐标集的并；投影满射保证两提升集合相同。有限乘积中添加闲置概率坐标不改变质量，所以 $P$ 良定义。在共同有限坐标集内计算又给有限可加性、非负性及 $P(X)=1$。

为应用 Carathéodory 延拓，尚须证明它是预测度。对有限可加概率，只需证明 $C_n\in\mathcal C$、$C_n\downarrow\varnothing$ 时 $P(C_n)\downarrow0$：如果不交 $A_j\in\mathcal C$ 的并 $A\in\mathcal C$，则 $A\setminus\bigcup_{j\leqslant n}A_j\downarrow\varnothing$，有限可加性便给可数可加性。

一个序列 $C_n$ 总共只使用至多可数个坐标。若这些坐标有限，在对应有限乘积空间用测度的从上连续性即可。若无限，列为 $1,2,\ldots$，可将 $C_n$ 写成依赖前 $k_n$ 个坐标的柱集，其中 $k_n$ 递增且 $k_n\geqslant n$。设反面情形 $P(C_n)\downarrow d>0$。对每个固定的前 $k$ 个坐标 $x_1,\ldots,x_k$，定义
\begin{align*}
p_n^k(x_1,\ldots,x_k)
=\int\chi_{C_n}(x_1,\ldots,x_k,x_{k+1},\ldots)
\,\mathrm d\mu_{k+1}\cdots\mathrm d\mu_{k_n}
\end{align*}
当 $k<k_n$；当 $k\geqslant k_n$ 时，直接取该柱集在已定坐标上的示性函数值。积分只有有限多个变量，未使用待构造的无限乘积。有限 Tonelli 定理保证这些函数可测、值在 $[0,1]$ 内。将不同 $n$ 提升到共同有限坐标集，$C_{n+1}\subseteq C_n$ 给逐点 $p_{n+1}^k\leqslant p_n^k$。令 $p^k=\lim_np_n^k$；有限乘积积分的相容性及有界收敛给出，对\emph{每个}前 $k$ 个坐标，
\begin{align*}
p^k(x_1,\ldots,x_k)
=\int p^{k+1}(x_1,\ldots,x_k,t)\,\mathrm d\mu_{k+1}(t).
\end{align*}
这里对固定前缀的可测截面应用有限 Tonelli，故等式不是仅对几乎所有前缀成立。零坐标值 $p^0=d>0$。递归地，若已选前缀使 $p^k>0$，上式说明存在 $x_{k+1}$ 使 $p^{k+1}>0$；由此选出一条所有前缀均正的序列。对固定 $n$，取 $k\geqslant k_n$，则 $0<p^k\leqslant p_n^k=\chi_{C_n}(x_1,\ldots,x_k)$，故所选点属于 $C_n$。对未参与的坐标任取一点，得到 $\bigcap_nC_n$ 中的点，矛盾。因此 $d=0$，预测度性成立。

Carathéodory 定理遂将 $P$ 延拓为 $\sigma(\mathcal C)=\bigotimes_{\alpha\in A}\mathscr M_\alpha$ 上的概率测度。原预测度有限，故延拓在这个 $\sigma$ 代数上唯一。这一构造给出全部有限维边缘恰为 $\mu_B$，特别给出有限矩形的乘积公式。

动机是把有限个坐标的独立概率读数扩展为一个完整测度。归一化为 $1$ 同时保证闲置坐标相容、有界收敛的统一界及延拓唯一性；没有假设各空间可度量、标准 Borel 或紧。可数操作只涉及可数坐标，递归选择则排除具有固定正概率的递减空柱集。选择任意坐标点沿用本书选择公理背景，因每个概率空间非空。空指标情形直接是单点概率空间。
\end{remark}
% END SUPPLEMENT: id=v14-ch07-notes-general-products-proof
% SOURCE p=248; printed=233; object=ch07-notes-product-domain-comparison; region=middle
\cref{theorem:ch07-product-regular-borel}的假设更强，但结论也更强；特别地，该定理中 $\mu$ 的定义域是整个乘积空间的 Borel $\sigma$ 代数。原书说，指标集 $A$ 不可数时，它比 $\bigotimes_\alpha\mathscr B_{X_\alpha}$ “大得多”。这个说法描述通常的非退化情形，若解释为对任意因子族都严格较大，则须作以下限定。
% SUPPLEMENT: id=v14-ch07-notes-product-domain-details; parent=ch07-notes-product-domain-comparison; kind=omitted-proof+source-clarification+motivation+hypotheses; origin=independent-derivation
\begin{remark}
\textbf{定义域严格扩大及退化边界（重构新增）。}任意乘积 $\sigma$ 代数中的集合都只依赖至多可数个坐标。事实上，所有具有这一性质的集合组成 $\sigma$ 代数：补集依赖同一组坐标，可数并依赖这些可数坐标集的可数并；它包含全部单坐标生成元，故包含整个乘积 $\sigma$ 代数。这里“只依赖 $J$”指两个点若在 $J$ 上相同，就同时属于或不属于该集合。

取不可数 $A$，每个 $X_\alpha=\{0,1\}$ 赋予离散拓扑，配上任意概率测度。这些因子紧 Hausdorff，测度 Radon。乘积中的单点 $Z=\{(0)_{\alpha\in A}\}$ 闭，所以 Borel；但它不只依赖任意给定的可数 $J$：选 $\beta\in A\setminus J$，把全零点的第 $\beta$ 坐标改成 $1$，得到 $J$ 上相同而不在 $Z$ 中的点。因此 $Z$ 不属于乘积 $\sigma$ 代数，给出严格包含的具体证明。

若所有因子均为单点，则不论 $A$ 多大，乘积仍为单点，两种 $\sigma$ 代数完全相同。因此“$A$ 不可数”本身不是严格包含的充分条件。拓扑版本的价值在于同时给出全 Borel 定义域与 Radon 正则性；这些信息不能由纯测度论乘积的定义域自动推出。
\end{remark}
% END SUPPLEMENT: id=v14-ch07-notes-product-domain-details
\end{document}
